% titel: Looping: Kugel rollt herab
% skript: Kap. Energiesatz, Rotation (Rollbewegung)

Eine Kugel (homogen, Trägheitsmoment $I_K=\frac25 m r^2$, Radius $r$) rollt
ohne zu gleiten aus einer Höhe $h$ in einen kreisrunden Looping (Radius $R$)
herab. Gesucht ist die Mindesthöhe $h$, damit die Kugel im höchsten Punkt
gerade nicht herunterfällt. Dabei ist neben der Translationsenergie auch die
Rotationsenergie zu berücksichtigen.

\begin{werkzeugkasten}[Werkzeuge]
\textbf{Energiesatz für rollende Kugel} (ohne Gleiten, ohne Reibungsverluste;
Rollbedingung $\omega=v/r$):
\begin{equation}
  m g \Delta h
  = \frac12 m v^2 + \frac12 I \omega^2
  = \frac12 m v^2 \left(1+\frac{r^2 I}{m r^2}\right)
  = \frac12 m v^2\left(1+\frac{I}{m r^2}\right).
\end{equation}

Für die Vollkugel $I=\frac25 m r^2$:
\begin{equation}
  m g\Delta h = \frac12 m v^2\left(1+\frac25\right) = \frac{7}{10} m v^2
  \;\Rightarrow\;
  v^2 = \frac{10}{7}\,g\,\Delta h.
\end{equation}

\textbf{Im höchsten Punkt des Loopings} (Radius $R$): Die Kugel darf gerade
nicht den Kontakt verlieren, d.h.\ die Normalkraft $\ge 0$; dann muss die
Zentripetalbeschleunigung ganz von der Schwerkraft geliefert werden:
\begin{equation}
  \frac{m v^2}{R} \ge m g
  \;\Rightarrow\;
  v^2 \ge g\, R.
\end{equation}
($R$ ist der Bahnradius des Kugelmittelpunkts.)

\textbf{Höhendifferenz} $\Delta h$: vom Loslasspunkt bis zum höchsten Punkt
des Loopings. In der Aufgabenstellung ist $h$ die \emph{zusätzliche} Höhe
über dem höchsten Punkt, also $\Delta h=h$.
\end{werkzeugkasten}

\begin{werkzeugkasten}[Vorgehen in allgemeiner Form]
Energiebilanz zwischen Loslassen (Ruhe) und höchstem Looping-Punkt
($\Delta h=h$); daraus $v^2=\frac{10}{7}gh$. Kontaktbedingung
$v^2\ge gR$ einsetzen und nach $h$ auflösen.
\end{werkzeugkasten}

\begin{achtung}[Typische Fallen]
\begin{itemize}
  \item Der $\frac25$-Anteil: Die Rotation "verbraucht" einen Teil der
       Energie; ohne ihn würde $h$ falsch klein herauskommen.
  \item Im Scheitelpunkt übernehmen \emph{Gewicht und Normalkraft} gemeinsam
       die Zentripetalbeschleunigung; "nicht herunterfallen" heisst
       $N\ge0$, also $v^2\ge gR$.
  \item Reihenfolge: erst $v^2(h)$, dann die Kontaktbedingung anwenden;
       nicht die Höhenangaben verwechseln.
\end{itemize}
\end{achtung}
